The evolution systems module contains the implementation of all PDE solvers in Thunder. Each system of equations is implemented as a lightweight class exposing a pre-\/defined interface that will be called from within the evolution loop in the code. At a high level, an evolution system holds the {\ttfamily state} array and the {\ttfamily aux} array as member variables and needs to implement the following functionality\+:
\begin{DoxyItemize}
\item {\bfseries{Initial data}}\+: The initial data is set right after initialization in Thunder, meaning that all the infrastructural data structures are initialized and in a valid state when this function is called. This includes all the grid data structures, such as the oct-\/tree forest, the connectivity, and the individual quadrants forming the simulation grid. When the initial data routine is called the first time, the grid is always uniformly spaced (in logical coordinates, i.\+e. all trees of the forest are at a uniform refinement level), with the level of refinement set by the {\itshape amr\+::initial\+\_\+refinement\+\_\+level} parameter. Optionally, this routine can be called a second time after the first round of grid refinement. This behaviour is controlled by the {\itshape evolution\+::reset\+\_\+id\+\_\+after\+\_\+regrid} parameter. Upon exiting this routine, it is expected that all variables in the {\ttfamily state} array (i.\+e. all the evolved variables) are in a valid state at {\bfseries{all}} cells, ghostzones included. The initial data is usually implemented as a free function outside of the evolution system class.
\item {\bfseries{Auxiliaries computation}}\+: Auxiliary variables pertaining to a certain evolution module need to be computed within the module itself. This is implemented as a {\ttfamily void} method called {\ttfamily compute\+\_\+auxiliaries} that takes the cell indices {\ttfamily i,j,k} and the quadrant index {\ttfamily q} as input. The method {\bfseries{must}} be marked {\ttfamily ALWAYS\+\_\+\+INLINE} as well as {\ttfamily HOST\+\_\+\+DEVICE} since this kernel will be called from within a parallel region. This function is expected to fill all the relevant entries in the {\ttfamily aux} array at {\bfseries{all}} cells, including ghost zones. This method has to be marked {\ttfamily const}.
\item {\bfseries{Maximum eigenspeed computation}}\+: In order to determine the largest stable timestep it is necessary to know the maximum eigenspeed associated to all the active evolution systems. This is implemented in the class as a method called {\ttfamily compute\+\_\+max\+\_\+eigenspeed} which again takes cells and quadrant indices as input. The method is expected to return the maximum absolute value eigenspeed as a {\ttfamily double}. This method also needs to be marked {\ttfamily ALWAYS\+\_\+\+INLINE} as well as {\ttfamily HOST\+\_\+\+DEVICE}.
\item {\bfseries{Flux computation}}\+: For a high-\/resolution shock capturing scheme the fluxes are at the heart of the evolution process. In Thunder, evolution is performed using dimensional flux-\/splitting, and for this reason an evolution system class has to implement three methods called {\ttfamily compute\+\_\+(x/y/z)\+\_\+flux} which fill the flux arrays. These methods are {\ttfamily void} and will be called from within a device kernel, so they must be tagged as {\ttfamily ALWAYS\+\_\+\+INLINE} as well as {\ttfamily HOST\+\_\+\+DEVICE} abd {\ttfamily const}. The methods take as input the cell indices {\ttfamily i,j,k} and the quadrant index {\ttfamily q}, as well as the number of ghostzones and the flux array to be filled. At this stage (i.\+e. during evolution), the internal state array contained by the evolution class is the correct one on which the fluxes need to be calculated. Reconstruction of variables happens within the flux computation kernel, and reconstructed variables are generally not stored by Thunder in order to optimize memory usage. The specific type of reconstruction used is normally a template parameter of the evolution class, and the reconstructor interface can be consulted in the documentation of {\ttfamily \mbox{\hyperlink{reconstruction_8hh_source}{reconstruction.\+hh}}}. The riemann solver is also a simple {\ttfamily struct} which is taken as a template parameter for the evolution system class. Each flux computation kernel will be called exactly on the indices where the fluxes need to be filled. The structure of the flux array is as follows\+: the first {\ttfamily NSPACEDIM} indices are cell indices, followed by the variable index (fluxes are only allocated for variables registered for hrsc evolution, see the documentation of the variable indices module for details). The next index in the array represents the direction of the flux, with {\ttfamily 0==X}, {\ttfamily 1==Y}, {\ttfamily 2==Z}, meaning that each kernel should fill the appropriate section of the array. Finally, the last index represents the quadrant index. The flux array does {\bfseries{not}} have ghost-\/zones allocated, meaning that the x-\/flux has dimensions {\ttfamily (nx+1) ny nz}, the y array has dimensions {\ttfamily nx (ny+1) nz} and the z array has dimensions {\ttfamily nx ny (nz+1)}. By convention, the x-\/flux entry at {\ttfamily (i,j,k)} represents the flux across the {\ttfamily (i-\/1/2,j,k)} face. This means that reconstruction routines provide reconstructed values consistent with this convention.
\item {\bfseries{Sources computation}}\+: Geometrical source terms are added in place in the new state array, and this is done within a method of the evolution class. The method must be called {\ttfamily compute\+\_\+source\+\_\+terms} and takes the cell and quadrant indices as well as the new state and timestep size as inputs. The method is expected to compute the source terms based on the internal state and add them in place onto the new state it takes as input.
\end{DoxyItemize}

For a more detailed description of an evolution system class, refer to the documentations of examples such as {\ttfamily burgers\+\_\+equation\+\_\+system\+\_\+t} or {\ttfamily scalar\+\_\+advection\+\_\+system\+\_\+t}. Each evolution system should define its own set of parameters. 